\section{Three convergent polylogarithmic integrals at every index}\label{tsc:polylog-moments}
The coefficient formula has a particularly explicit integral counterpart.
Let
\begin{equation}\label{tsc:cdalpha}
 \ell=\log(2\pi),\qquad c=\ell+\frac{\pi\ii}{2},\qquad
 d=\ell-\frac{\pi\ii}{2},\qquad \alpha=\gamma+c.
\end{equation}
For a nontrivial unit color $z$ and $x\ge0$, define
\begin{equation}\label{tsc:Vdef}
 V_m(z,x)=m![u^{m+1}]\,
      \Gamma(1+u)e^{-cu}\Li_u(ze^{-x}).
\end{equation}
This is a finite expression in ordinary polylogarithm order derivatives.
If
\[
 E_+(u)=\Gamma(1+u)e^{-cu}=\sum_{j\ge0}e_j^+u^j,
\]
then
\begin{equation}\label{tsc:Vfinite}
 V_m(z,x)=m!\sum_{r=0}^{m+1}\frac{e_{m+1-r}^+}{r!}
                         \Li_0^{[r]}(ze^{-x}),
\end{equation}
where
\begin{equation}\label{tsc:e-recurrence}
 e_0^+=1,\qquad
 n e_n^+=-\alpha e_{n-1}^+
          +\sum_{k=2}^n(-1)^k\zeta(k)e_{n-k}^+.
\end{equation}
For example, with $Z_r=\Li_0^{[r]}(ze^{-x})$,
\begin{align}\label{tsc:Vexamples}
 V_0&=Z_1-\alpha Z_0,\notag\\
 V_1&=\frac{Z_2}{2}-\alpha Z_1+
                  \frac{\alpha^2+\zeta(2)}2Z_0,\notag\\
 V_2&=\frac{Z_3}{3}-\alpha Z_2+
      (\alpha^2+\zeta(2))Z_1
      -\frac{\alpha^3+3\alpha\zeta(2)+2\zeta(3)}3 Z_0.
\end{align}
In particular, $Z_0(q)=q/(1-q)$ is elementary.

For $F$ analytic at zero and exponentially decreasing at infinity, put
\begin{align}\label{tsc:Mfunctional}
 \calM_{n,d}[F]={}&
    \int_0^1\frac{F(x)-F(0)}x(\log x-d)^n\dd x
   +\int_1^\infty\frac{F(x)}x(\log x-d)^n\dd x\notag\\
 &+\frac{(-d)^{n+1}}{n+1}F(0).
\end{align}
Both displayed integrals are ordinary absolutely convergent integrals.
The last term is essential; it is the contribution of the simple Mellin
pole after multiplication by $e^{-du}$.

\begin{theorem}[All-index convergent Mellin formula]\label{tsc:mellin-all}
With the colors \eqref{tsc:colors} and the conventions above,
\begin{equation}\label{tsc:M-all}
 \boxed{\displaystyle
 \calJ_{\mathbf m}(\mathbf a)=2\Re\sum_{k=1}^3
 \calM_{m_k,d}\!\big[
       V_{m_i}(z_{ik},\cdot)V_{m_j}(z_{jk},\cdot)\big].}
\end{equation}
Thus every separated triple Stieltjes finite part has a representation by
three convergent logarithmic integrals and three explicit endpoint terms.
Only finitely many polylogarithm order derivatives at order zero occur.
\end{theorem}
\begin{proof}
Consider the positive-frequency sector with singleton $k$. Before Laurent
coefficient extraction, multiply its Tornheim term by the three Gamma
factors appearing in \eqref{tsc:Zlaurent}. Inserting
\eqref{tsc:T-mellin} cancels the Gamma factor belonging to the singleton.
The sector becomes the analytic continuation of
\begin{equation}\label{tsc:sector-mellin}
 \Gamma(u_i)e^{-cu_i}\Gamma(u_j)e^{-cu_j}e^{-du_k}
 \int_0^\infty x^{u_k-1}
           \Li_{u_i}(z_{ik}e^{-x})\Li_{u_j}(z_{jk}e^{-x})\dd x.
\end{equation}
Taking the nonnegative Laurent coefficients in $u_i,u_j$, with the
factorials prescribed by \eqref{tsc:stieltjes-convention}, replaces the two
polylogarithmic factors by $V_{m_i}$ and $V_{m_j}$.

For the remaining variable $u=u_k$, write the Mellin integral as
\[
 \frac{F(0)}u+
 \int_0^1x^{u-1}(F(x)-F(0))\dd x+
 \int_1^\infty x^{u-1}F(x)\dd x.
\]
The coefficient $m_k![u^{m_k}]$ after multiplication by $e^{-du}$ is exactly
\eqref{tsc:Mfunctional}. The nontrivial colors ensure $F(x)-F(0)=O(x)$;
all of its spectral derivatives have the same property. At infinity they
decay exponentially. These local bounds justify every finite spectral
extraction under the subtracted integrals. Reversing the three Fourier
signs conjugates this expression because the shifts and Stieltjes indices
are real. Summing the three positive sectors and their conjugates proves
\eqref{tsc:M-all}.
\end{proof}

\begin{corollary}[A convergent three-digamma formula]\label{tsc:psi-mellin}
Put
\[
 W_z(x)=\left.\partial_s\Li_s(ze^{-x})\right|_{s=0}
            -\left(A+\frac{\pi\ii}{2}\right)\frac{ze^{-x}}{1-ze^{-x}}.
\]
Then
\begin{align}\label{tsc:psi-integrals}
 \FP\int_0^1\prod_{i=1}^3\psi(\{x-a_i\})\dd x
 =-2\Re\sum_{k=1}^3\bigg\{&
  \int_0^1\frac{W_{z_{ik}}(x)W_{z_{jk}}(x)
                 -W_{z_{ik}}(0)W_{z_{jk}}(0)}x\dd x\notag\\
 &+\int_1^\infty\frac{W_{z_{ik}}(x)W_{z_{jk}}(x)}x\dd x
   -dW_{z_{ik}}(0)W_{z_{jk}}(0)\bigg\}.
\end{align}
\end{corollary}
All integrals on the right are ordinary convergent integrals. For example,
at $(a_1,a_2,a_3)=(0,1/4,2/3)$ the common value is approximately
\[
 30.589007126644942529054519724157914\ldots.
\]
This decimal is a diagnostic, not a claim of a reduction to a smaller
constant basis. The accompanying code evaluates the left side by its
three local endpoint subtractions and the right side by convergent
polylogarithmic series; their residual at 40 working digits is below
$3\cdot10^{-38}$.

\subsection{Rational colors make the endpoint data finite}
Let $z=e(p/q)\ne1$ with integers $0<p<q$. Splitting the polylogarithm
series into residue classes gives the entire identity
\begin{equation}\label{tsc:finite-DFT}
 \Li_u(z)=q^{-u}\sum_{r=1}^q z^r\zeta(u,r/q).
\end{equation}
The apparent pole on the right cancels because $\sum_{r=1}^qz^r=0$.
Lerch's formula $\zeta^{[1]}(0,x)=\log\Gamma(x)-\ell/2$ then gives
\begin{equation}\label{tsc:W-boundary}
 W_z(0)=\sum_{r=1}^qz^r\log\Gamma(r/q)
       -(\alpha+\log q)\frac{z}{1-z}.
\end{equation}
More generally, substituting \eqref{tsc:finite-DFT} into \eqref{tsc:Vdef}
gives every $V_m(z,0)$ as a finite combination of rational-argument
Hurwitz jets at zero. Thus the endpoint terms in the rational-shift
three-digamma formula are explicitly Gamma-valued, even when the remaining
convergent integrals have not been reduced further.
